\documentclass[10pt]{article}
\usepackage[T1]{fontenc}
\usepackage[utf8]{inputenc}
\usepackage[french]{babel}
\usepackage{amsmath,amssymb,mathrsfs,mathpazo}
\usepackage{geometry}
\geometry{a4paper,margin=2cm}
\usepackage{graphicx,array,tabularx,tabularray,colortbl,multirow,multicol,booktabs,enumitem,pifont,ulem,fancybox,diagbox,cancel}
\usepackage{pstricks,pstricks-add,pst-plot,pst-tree,pst-node,pst-text,pst-func,pst-math,pst-bspline,pst-3dplot,pst-eucl,pst-solides3d}
\let\pscalculate\undefined
\usepackage{pst-bezier}
\usepackage{tikz,pgfplots,tkz-tab}
\usetikzlibrary{arrows,arrows.meta,patterns,calc,positioning,trees,decorations.markings,shapes.geometric,shapes.misc,shapes.arrows,intersections,decorations.pathmorphing,decorations.pathreplacing,backgrounds,shadows,babel}
\pgfplotsset{compat=1.18}
\usepgfplotslibrary{fillbetween}
\usepackage[np]{numprint}
\usepackage{listings,makecell,pas-tableur,twoopt}
\let\barre\undefined
\usepackage[pstricks]{tablvar}
\definecolor{darkgreen}{RGB}{0,100,0}
\providecommand{\red}{\color{red}}
\providecommand{\blue}{\color{blue}}
\providecommand{\green}{\color{green}}
\providecommand{\black}{\color{black}}
\providecommand{\white}{\color{white}}
\providecommand{\gray}{\color{gray}}
\providecommand{\cyan}{\color{cyan}}
\providecommand{\magenta}{\color{magenta}}
\providecommand{\yellow}{\color{yellow}}
\providecommand{\violet}{\color{violet}}
\providecommand{\orange}{\color{orange}}
\providecommand{\euro}{\texteuro}
\providecommand{\EUR}{\texteuro}
\providecommand{\ieme}{\textsuperscript{e}}
\providecommand{\degres}{\ensuremath{^\circ}}
\providecommand{\degree}{\ensuremath{^\circ}}
\providecommand{\vv}[1]{\overrightarrow{#1}}
\providecommand{\Coord}[1]{#1}
\providecommand{\decosix}{}
\providecommand{\starredbullet}{\ensuremath{\bullet}}
\pagestyle{empty}
\setlength{\parindent}{0pt}
\usepackage{adjustbox}\DeclareUnicodeCharacter{2236}{:}\DeclareUnicodeCharacter{2190}{\ensuremath{\leftarrow}}
\providecommand{\R}{\mathbb{R}}
\providecommand{\N}{\mathbb{N}}
\providecommand{\Z}{\mathbb{Z}}
\providecommand{\Q}{\mathbb{Q}}
\providecommand{\D}{\mathbb{D}}
\providecommand{\C}{\mathbb{C}}
\providecommand{\e}{{\mathrm{e}}}
\providecommand{\vect}[1]{\overrightarrow{#1}}
\providecommand{\vectt}[1]{\overrightarrow{\mathrm{#1}}}
\providecommand{\vvt}[1]{\overrightarrow{\mathrm{#1}}}
\providecommand{\np}[2][]{{#2~#1}}
\providecommand{\vv}[1]{{\overrightarrow{#1}}}
\renewcommand{\d}{{\mathrm{d}}}
\providecommand{\degres}{{^\circ}}
\providecommand{\degre}{{^\circ}}
\providecommand{\blue}{\color{blue}}
\providecommand{\red}{\color{red}}
\providecommand{\enskip}{\quad}
\providecommand{\llbracket}{\lbrack\!\lbrack}
\providecommand{\rrbracket}{\rbrack\!\rbrack}
\providecommand{\euro}{{\text{€}}}
\providecommand{\no}{{\text{n°}}}
\providecommand{\interval}[2]{{[#1;#2]}}
\providecommand{\intervFO}[2]{{[#1;#2[}}
\providecommand{\dotfill}{\ldots}

\providecommand{\ds}{\displaystyle}
\providecommand{\textding}[1]{#1}
\providecommand{\Oij}{\(\mathrm{O};\vec{\imath},\vec{\jmath}\)}
\providecommand{\Oijk}{\(\mathrm{O};\vec{\imath},\vec{\jmath},\vec{k}\)}
\providecommand{\Ouv}{\(\mathrm{O};\vec{u},\vec{v}\)}
\providecommand{\intervFF}[2]{\ensuremath{[#1;#2]}}
\renewcommand{\euro}{\ifmmode\text{\texteuro}\else\texteuro\fi}
\DeclareUnicodeCharacter{00B0}{\ifmmode^{\circ}\else\textdegree\fi}
\DeclareUnicodeCharacter{20AC}{\ifmmode\text{\texteuro}\else\texteuro\fi}
% Énoncés : Éducation nationale. Transcription et corrigés : APMEP, Denis Vergès et contributeurs. Découpage et classement : CoopMaths, Sébastien Lozano, Éric Elter et contributeurs. Adaptation de la présentation : Maths Explorer.
% CC BY-SA : https://coopmaths.fr/www/about/
% https://coopmaths.fr/alea/static/bac/2022/tex/bac_2022_10_sujet2_caledonie_2_cor.tex
\begin{document}
\noindent\textit{APMEP / CoopMaths. CC BY-SA. Présentation adaptée par Maths Explorer.}\par\medskip
\textbf{Principaux domaines abordés :} fonctions, fonction logarithme.

\medskip

Soit $f$ la fonction définie sur l'intervalle $]0~;~+\infty[$ par :

\[f(x) = x \ln (x) - x - 2.\]







\begin{enumerate}
\item 
	\begin{enumerate}
		\item 
		Somme de produits de fonctions dérivables sur $]0~;~+\infty[$, \: $f$ est dérivable et sur cet intervalle :
		
		$f'(x) = \ln x + x \times \dfrac1x - 1 = \ln x + 1 - 1 = \ln x$.
		\item 
On a $f(\text{e}) = \e \ln \e - \e - 2 = - 2$.

De plus $f'(\e) = \ln \e = 1$.

On sait que $M(x~;~y) \in (T) \iff y - (- 2) = 1(x - \e) \iff y = x - 2 - \e$.
		\item 
De $f'(x) = \ln x$, on en déduit que $f''(x) = \dfrac1x$.

Comme $x > 0$, on a donc $f''(x) > 0$ sur $]0~;~+\infty[$ : la fonction $f$ est donc convexe sur cet intervalle.
		\item 
		Le résulta précédent montre que la courbe $\mathcal{C}_f$ est au dessus de la tangente $T$.
	\end{enumerate}	
\item 
	\begin{enumerate}
		\item 
On sait que $\displaystyle\lim_{x \to 0} x \ln x = 0$, que $\displaystyle\lim_{x \to 0} x  = 0$, donc par somme de limites $\displaystyle\lim_{x \to 0} f(x) = - 2$
		\item 
On a pour tout réel $x \in ]0~;~+ \infty[$, on peut écrire :

$f(x) = x\left[\ln (x) - 1\right] - 2$.

Comme $\displaystyle\lim_{x \to + \infty} \ln x - 1 = + \infty$ et $\displaystyle\lim_{x \to + \infty} x = + \infty$, on a par produit de limites : $\displaystyle\lim_{x \to + \infty} f(x) = + \infty$.
	\end{enumerate}
\item 
On a $f'(x) = \ln x$. On sait que :

$\bullet~~$sur $]0~;~1[, \: \ln x < 0$ : donc $f$ décroit sur cet intervalle ;

$\bullet~~$sur $]1~;~+\infty[, \: \ln x > 0$ : donc $f$ croit sur cet intervalle ;

$\bullet~~$$f(1) = - 1 - 2 = - 3$ est donc le minimum de $f$ sur $]0~;~+ \infty[$.

D'où le tableau de variations :

\begin{center}
\psset{unit=1cm,arrowsize=2pt 3}
\begin{pspicture}(7,3)
\psframe(7,3)\psline(0,2)(7,2)\psline(0,2.5)(7,2.5)\psline(1,0)(1,3)
\psline(1.15,0)(1.15,2.5)\psline(1.2,0)(1.2,2.5)
\uput[u](0.5,2.4){$x$}\uput[u](1.15,2.4){$0$}\uput[u](4,2.4){$1$}\uput[u](6.5,2.4){$+\infty$}
\rput(0.5,2.25){$f'(x)$}\rput(2.5,2.25){$-$}\rput(4,2.25){0}\rput(5.5,2.25){$+$}
\uput[d](1.4,2){$-2$}\uput[u](4,0){$-3$}\uput[d](6.5,2){$+ \infty$}
\psline{->}(1.5,1.5)(3.5,0.5)\psline{->}(4.5,0.5)(6.5,1.5)
\rput(0.5,1){$f$}
\end{pspicture}

\end{center}
\item
	\begin{enumerate}
		\item 

Le tableau de variations montre que sur l'intervalle $]1~;~+ \infty[$, $f$ est strictement croissante de $- 3$ à plus l'infini.

$f$ étant continue car dérivable, le théorème des valeurs intermédiaires montre qu'il existe un réel unique $\alpha$ de $]1~;~+ \infty[$ tel que $f(\alpha) = 0$.
		\item 
		On a $f(4,3) \approx - 0,028$ et $f(4,4) \approx 0,119$.
		
D'après le même théorème ceci montre que $\alpha \in ]4,3~;~4,4[$.
		\item 
		Conclusion :
		
$\bullet~~$sur $]0~;~\alpha[, \: f(x) < 0$ ;

$\bullet~~$sur $]\alpha~;~+ \infty[, \: f(x) > 0$ ;

$\bullet~~$$f(\alpha) = 0$.
	\end{enumerate}
\item 















La fonction \texttt{seuil(0.01)} renvoie la valeur 4,32.


Ceci donne l'encadrement de $\alpha$ au centième près :
$4,31 < \alpha < 4,32$.
\end{enumerate}

\bigskip

\end{document}
